%======================================================================
% CAPÍTULO 9: SANDRA CONCILIATION FILE (SCF)
%======================================================================
\chapter{Sandra Conciliation File: Sistema de Conciliación de Archivos}

\textbf{Sandra Conciliation File (SCF)} es una herramienta de comparación de archivos CSV de alto rendimiento, diseñada para identificar coincidencias y hallazgos entre dos fuentes de datos. Originalmente desarrollada como proyecto independiente, fue integrada en \textbf{Sentinel} como módulo de conciliación bancaria.

\notacientifica{SCF implementa un algoritmo de Hash Join en memoria, permitiendo сравнения eficiente de grandes archivos sin necesidad de bases de datos.}

\section{Introducción y Propósito}

La conciliación de archivos es un proceso fundamental en entornos empresariales donde es necesario comparar registros de diferentes fuentes: sistemas internos vs. extractos bancarios, nóminas vs. transferencias, o bases de datos legacy vs. sistemas modernos.

\begin{caja}{Caso de Uso}
\begin{itemize}
\item \textbf{Conciliación Bancaria}: Comparar movimientos del banco con registros del sistema de nómina.
\item \textbf{Auditoría de Datos}: Identificar discrepancias entre dos bases de datos.
\item \textbf{Validación de Carga}: Verificar que los datos cargados coinciden con los originales.
\item \textbf{Detección de Duplicados}: Encontrar registros duplicados entre sistemas.
\end{itemize}
\end{caja}

\notacientifica{La comparación eficiencia de archivos es crítica en entornos con miles o millones de registros, donde herramientas tradicionales como diff o grep resultan insuficientes.}

\subsection{Definición del Sistema}

\begin{definicion}
\textbf{SCF (Sandra Conciliation File)} es una herramienta CLI que compara dos archivos CSV delimitados, identificando:
\begin{itemize}
\item \textbf{Coincidencias}: Registros del origen que existen en el archivo de comparación.
\item \textbf{Hallazgos}: Registros del origen que NO existen en el archivo de comparación.
\item \textbf{Errores}: Problemas de formato en los archivos de entrada.
\end{itemize}
\end{definicion}

\notacientifica{SCF utiliza claves compuestas para la comparación, permitiendo definir qué columnas constituyen la identidad única de cada registro.}

\section{Arquitectura}

SCF está diseñado bajo principios de modularidad y eficiencia, con una arquitectura que facilita tanto el uso independiente como la integración en Sentinel:

\subsection{Estructura de Módulos}

\begin{caja}{Componentes del Sistema}
\begin{itemize}
\item \textbf{core/keys.rs}: Generador de claves de comparación.
\item \textbf{core/comparator.rs}: Motor de comparación Hash Join.
\item \textbf{io/mod.rs}: Lectura/escritura de archivos con buffering.
\item \textbf{report/mod.rs}: Generación de reportes de resultados.
\item \textbf{error.rs}: Manejo de errores personalizado.
\item \textbf{cli.rs}: Interfaz de línea de comandos.
\end{itemize}
\end{caja}

\notacientifica{La arquitectura modular permite reutilizar los componentes core en diferentes contextos, incluyendo la integración directa en Sentinel.}

\subsection{Patrón de Comparación (Hash Join)}

SCF implementa el algoritmo \textbf{Hash Join} para comparar archivos de manera eficiente:

\begin{caja}{Fases del Algoritmo}
\begin{enumerate}
\item \textbf{Carga (Build Phase)}: Lee el archivo de comparación y genera un HashSet de claves.
\item \textbf{Comparación (Probe Phase)}: Lee el archivo de origen y para cada registro verifica si su clave existe en el HashSet.
\item \textbf{Resultado}: Escribe los registros en archivos separados según el resultado.
\end{enumerate}
\end{caja}

\notacientifica{La complejidad del algoritmo Hash Join es O(N+M), significativamente más eficiente que algoritmos de comparación línea a línea.}

\bigskip

\subsection{Stack Tecnológico}

\begin{caja}{Tecnologías del Stack}
\begin{itemize}
\item \textbf{Lenguaje}: Rust (Edición 2021).
\item \textbf{CLI}: Clap para parsing de argumentos.
\item \textbf{Concurrencia}: Rayon para procesamiento paralelo.
\item \textbf{UI}: Indicatif para barras de progreso.
\item \textbf{Errores}: Thiserror para manejo de errores tipado.
\end{itemize}
\end{caja}

\notacientifica{El uso de Rayon permite paralelizar la comparación de grandes archivos, aprovechando todos los cores disponibles.}

\section{Uso del Sistema}

SCF puede utilizarse de dos maneras: como herramienta CLI independiente o como módulo integrado en Sentinel.

\subsection{Como Herramienta Independiente}

\begin{caja}{Sintaxis CLI}
\begin{verbatim}
scf <archivo_comparacion> -c <columnas> <archivo_origen> -o <columnas> [opciones]

Opciones:
  -d, --delimiter <CAR>    Delimitador (default: ;)
  -O, --output <DIR>       Directorio de salida (default: out)
  --skip-header            Omitir primera línea
  -q, --quiet              Modo silencioso
  -v, --verbose            Salida detallada
\end{verbatim}
\end{caja}

\notacientifica{La flexibilidad en la definición de claves permite adaptar la comparación a cualquier formato de archivo.}

\bigskip

\subsection{Ejemplo de Uso}

\begin{caja}{Ejemplo Práctico}
\begin{verbatim}
# Comparar movimientos bancarios vs. sistema de nómina
scf banco.csv -c 0,1,2 nomina.csv -o 0,1,3

# Output:
# out/
# ├── coincidencias.txt   # Registros que coinciden
# ├── hallazgos.txt       # Registros en banco no en sistema
# ├── log.txt           # Errores encontrados
# └── reporte.txt      # Resumen de la operación
\end{verbatim}
\end{caja}

\notacientifica{La estructura de salida facilita la integración con flujos de trabajo de auditoría.}

\subsection{Como Comando de Sentinel}

SCF está integrado en Sentinel como el comando \texttt{conciliate}:

\begin{caja}{Sintaxis en Sentinel}
\begin{verbatim}
sandra-sentinel conciliate \
  --comparison <archivo> \
  --comparison-columns <col1,col2,...> \
  --origin <archivo> \
  --origin-columns <col1,col2,...> \
  [--delimiter <char>] \
  [--output <dir>] \
  [--skip-header] \
  [--quiet]
\end{verbatim}
\end{caja}

\notacientifica{La integración en Sentinel permite combinar la conciliación de archivos con el resto de capacidades del motor de cálculo.}

\section{Optimizaciones de Rendimiento}

SCF implementa varias optimizaciones para manejar archivos de gran tamaño:

\subsection{Pre-asignación de Memoria}

\begin{caja}{Optimización de HashSet}
\begin{itemize}
\item Uso de \texttt{HashSet::with\_capacity(capacity)} para evitar rehashings.
\item Pre-cálculo del tamaño basado en el archivo de comparación.
\item Reducción significativa de asignaciones de memoria.
\end{itemize}
\end{caja}

\notacientifica{La pre-asignación puede reducir el tiempo de ejecución en hasta 30\% para archivos grandes.}

\bigskip

\subsection{Buffered I/O}

\begin{caja}{E/S con Buffer}
\begin{itemize}
\item \texttt{BufReader} con capacidad de 1MB para lectura.
\item \texttt{BufWriter} para escritura eficiente.
\item Reducción del número de syscalls al sistema operativo.
\end{itemize}
\end{caja}

\notacientifica{El buffering es especialmente efectivo en sistemas con latency de disco alta.}

\bigskip

\subsection{Procesamiento Paralelo}

\begin{caja}{Paralelización con Rayon}
\begin{itemize}
\item División del archivo de origen en chunks de 1000 líneas.
\item Procesamiento paralelo de cada chunk.
\item Reconstrucción de resultados en orden.
\end{itemize}
\end{caja}

\notacientifica{Rayon transparenta la paralelización, permitiendo escalar a múltiples cores sin cambiar la lógica del código.}

\bigskip

\subsection{Barra de Progreso}

\begin{caja}{Indicatif Progress Bar}
\begin{itemize}
\item Visualización en tiempo real del progreso.
\item Estimación de tiempo restante.
\item Actualización cada 10,000 líneas.
\end{itemize}
\end{caja}

\notacientifica{La barra de progreso es especialmente útil para archivos de millones de líneas donde la operación puede tomar varios minutos.}

\section{Formato de Archivos}

SCF soporta archivos CSV con delimitador configurable:

\begin{caja}{Especificación de Formato}
\begin{itemize}
\item \textbf{Delimitador}: Por defecto es punto y coma (;), configurable.
\item \textbf{Codificación}: UTF-8.
\item \textbf{Salto de línea}: LF o CRLF.
\item \textbf{Encabezados}: Opcional, con flag \texttt{--skip-header}.
\end{itemize}
\end{caja}

\notacientifica{La flexibilidad en el delimitador permite trabajar con diferentes estándares regionales.}

\subsection{Definición de Claves}

Las claves de comparación se definen mediante índices de columnas (0-based):

\begin{caja}{Ejemplos de Claves}
\begin{verbatim}
# Clave simple (una columna)
-c 0

# Clave compuesta (múltiples columnas)
-c 0,1,2

# Diferentes claves por archivo
-c 0,1,2  (comparación usa columnas 0,1,2)
-o 0,3,4  (origen usa columnas 0,3,4)
\end{verbatim}
\end{caja}

\notacientifica{Las claves compuestas permiten comparación precisa cuando un solo campo no es único.}

\section{Salida del Sistema}

\subsection{Archivos Generados}

\begin{caja}{Archivos de Resultado}
\begin{enumerate}
\item \textbf{coincidencias.txt}: Registros del origen que existen en la comparación.
\item \textbf{hallazgos.txt}: Registros del origen que NO existen en la comparación.
\item \textbf{log.txt}: Errores de formato encontrados.
\item \textbf{reporte.txt}: Resumen con conteo de líneas y tamaño de archivos.
\end{enumerate}
\end{caja}

\notacientifica{La separación de resultados facilita el flujo de trabajo de auditoría.}

\subsection{Formato del Reporte}

\begin{caja}{Estructura del Reporte}
\begin{verbatim}
----------------------------------------
Archivo          Líneas   Peso (bytes)
----------------------------------------
coincidencias.txt   1523   89234
hallazgos.txt        287   12456
log.txt                0      0
\end{verbatim}
\end{caja}

\notacientifica{El reporte permite evaluar rápidamente el resultado de la conciliación.}

\section{Integración con Sentinel}

La integración de SCF en Sentinel permite un flujo de trabajo unificado:

\begin{caja}{Flujo de Trabajo Integrado}
\begin{enumerate}
\item Sentinel carga datos de múltiples fuentes (Banco, Sistema, Legacy).
\item SCF realiza la conciliación entre las fuentes.
\item El resultado se pasa al Motor de Cálculo para procesamiento.
\item Se genera la nómina auditada y validada.
\end{enumerate}
\end{caja}

\notacientifica{La integración elimina la necesidad de herramientas externas para la conciliación.}

\section{Conclusiones del Capítulo}

Sandra Conciliation File representa una herramienta especializada dentro del ecosistema Sandra:

\begin{caja}{Resumen de Capacidades}
\begin{itemize}
\item \textbf{Comparación Eficiente}: Algoritmo Hash Join O(N+M).
\item \textbf{Formato Flexible}: Delimitador configurable.
\item \textbf{Claves Compuestas}: Múltiples columnas para identificación.
\item \textbf{Rendimiento Optimizado}: Buffered I/O, pre-alloc, paralelo.
\item \textbf{Salida Estructurada}: Archivos separados para auditoría.
\item \textbf{Integración Total}: Comando nativo en Sentinel.
\end{itemize}
\end{caja}

\notacientifica{SCF complementa las capacidades de Sentinel, proporcionando una solución completa para la gestión y validación de datos en entornos empresariales.}
